% ==============================================================================
% Fichier : front/02_mode_emploi.tex : SGSI (SIS498)
% ==============================================================================

\chapter*{Mode d'emploi}
\addcontentsline{toc}{chapter}{Mode d'emploi}
\markboth{Mode d'emploi}{Mode d'emploi}

\section*{À qui s'adresse ce cours}

Aux étudiants du Master Professionnel en Sécurité des Systèmes Informatiques
et des Données, et aux professionnels qui exercent ou se destinent aux
métiers de \textbf{RSSI}, de \textbf{DSI}, de \textbf{délégué à la protection
des données} (\sigle{DPO}), d'\textbf{auditeur de SMSI} et de
\textbf{responsable conformité}. Des notions de réseaux, de systèmes et de
droit numérique sont utiles.

\section*{Comment lire ce livre}

Le livre comprend un prologue et six parties. Chaque chapitre suit le même
plan~: une épigraphe, le niveau et les prérequis, les objectifs, une question
directrice, cinq sections, un encadré « à retenir », dix exercices gradués et
une transition. Les niveaux des exercices sont signalés par des étoiles~:
\niveaubase{} fondamental, \niveaumoyen{} intermédiaire, \niveauexpert{}
expert.

\section*{Les encadrés}

\begin{itemize}
    \item \textbf{Scénario opérationnel}~: un épisode des opérations NYONG,
          KELLÉ, MAKAK ou MATOMB, avec \emph{constatation} (faits) et
          \emph{analyse} (interprétation, niveau de confiance).
    \item \textbf{Actif}~: la fiche d'un actif informationnel (valeur, support,
          critères de sécurité prioritaires).
    \item \textbf{Mesure}~: une mesure de sécurité, son objectif et sa mise en
          \oe uvre.
    \item \textbf{Application métier}~: ce que le chapitre change pour un
          métier.
    \item \textbf{Cadre juridique}~: le texte de droit applicable.
\end{itemize}

\section*{Constatation et analyse}

Le cours distingue systématiquement ce que l'on \textbf{constate} (un fait
vérifiable, une preuve) de ce que l'on en \textbf{déduit} (une interprétation,
assortie d'un niveau de confiance). C'est la discipline première du management
de la sécurité et de l'audit.

\section*{NDAP SUKULU et évaluation}

L'environnement compagnon NDAP SUKULU (page suivante) propose, pour chaque
chapitre, une fiche d'auto-évaluation et, pour les chapitres pratiques, un
lab sur pièces. Fiches et labs se font hors des 30~heures de cours. Le code
du cours est \textbf{\codeUE}.
